\section{System Design}\label{sec:design}
% Source of truth: docs/ARCHITECTURE.md, docs/THREAT_MODEL.md and src/. Describe in your own words.

\subsection{Architecture}
\note{A shared \texttt{CredentialRegistry} plus stateless \texttt{IApprovalPolicy} implementations. The policy is
\texttt{immutable}, so it cannot be downgraded. Only issuance authorization differs between policies; the
credential format, storage, lifecycle, revocation and events are identical. Consider a small architecture
figure (draw it yourselves, e.g.\ in TikZ).}

\subsection{Actors and roles}
\note{Governor (manages universities and approvers, has no issuance power, cannot be an approver); approvers
with exactly one role each (Registrar, Dean, Exam Officer); verifier; holder. Static separation of duty.}

\subsection{Credential commitment}
\note{Give the exact procedure:
$c = \mathrm{keccak256}(\mathrm{abi.encode}(\mathit{DOMAIN}, \mathrm{keccak256}(\mathit{canonicalJSON}), \mathit{salt}))$,
with a 32-byte random salt held by the graduate and
$\mathit{DOMAIN}=\mathrm{keccak256}(\texttt{KZ-ACADEMIC-CREDENTIAL-v1})$.
Explain why the salt is needed (low-entropy fields) and why the verifier must cross-check the issuer named in
the document. Mention the Python--Solidity parity test over 50 synthetic credentials.}

\subsection{Lifecycle and verification states}
\note{None $\rightarrow$ Pending $\rightarrow$ Issued $\rightarrow$ Revoked; Cancelled burns the commitment.
The verifier sees NotFound, Pending, Valid or Revoked, plus the issuer and an issuer-active flag.}

\subsection{Approval policies}
\begin{table}[t]
\centering\small
\caption{Approval policies compared (university with 2 Registrars, 3 Deans and 2 Exam Officers).}
\label{tab:policies}
\begin{tabular}{lll}
\toprule
Policy & Eligible roles & Rule at finalization\\
\midrule
A Single issuer & Registrar & $\geq 1$ valid approval\\
B Unanimous & all & all $N=7$ active approvers\\
C Threshold & all & $\geq 4$ valid approvals\\
D Role-constrained & all & $\geq 4$ and $\geq 1$ per role\\
\bottomrule
\end{tabular}
\end{table}
\note{Explain why C and D share the threshold: it isolates the effect of role separation. Explain the
approver-nonce rule: any role change or (de)activation invalidates that approver's pending approvals, and
validity is re-checked at finalization.}

\subsection{Threat model}
\note{In-scope adversaries: a single insider, a $k$-coalition, an outsider, a document tamperer. Trusted: the
governor. Out of scope: falsified source records. List each enforced property with the test that checks it
(unauthorized issuance, duplicate issuance or approval, approval reuse, administrative bypass).}
